%% Chapter 10, Section 10.2 — Expectation Maximization Algorithm
%% Source: A First Course in Monte Carlo Methods, Sanz-Alonso & Al-Ghattas
%% Book pages 119–122

\section{Expectation Maximization Algorithm}
\label{sec:10.2}

In its most basic form, the Expectation Maximization (EM) algorithm is a deterministic optimization
method. EM is widely used to compute maximum likelihood estimators in problems where introducing a
latent variable makes the likelihood function easier to optimize.

Suppose for concreteness that we want to find the maximum likelihood estimator
\[
  \hat{\theta} = \argmax_{\theta}\, f(y|\theta)
\]
of a parameter $\theta$ given observed data $y$. Suppose that $f(y|\theta)$ is the marginal of a
complete likelihood $f(y, z|\theta)$, which includes the data $y$ and an auxiliary variable $z$.
Then, for any p.d.f.\ $g$ with compatible support, it holds by Jensen's inequality that
\begin{align*}
  \log f(y|\theta)
    &= \log \int f(y, z|\theta)\,dz \\
    &= \log \int \frac{f(y,z|\theta)}{g(z)}\, g(z)\,dz \\
    &\geq \int \log\!\left(\frac{f(y,z|\theta)}{g(z)}\right) g(z)\,dz \\
    &= \mathbb{E}_g[\log f(y, z|\theta)] - \mathbb{E}_g[\log g] = \ELBO(g, \theta).
\end{align*}
Here we view the ELBO as a function of the auxiliary distribution $g$ and the parameter $\theta$ in
the density $f(y, z|\theta)$. Now, notice that we have
\begin{align*}
  \ELBO(g, \theta)
    &= \int \log\!\left(\frac{f(y,z|\theta)}{g(z)}\right) g(z)\,dz \\
    &= \int \log\!\left(\frac{f(z|y,\theta)}{g(z)}\right) g(z)\,dz + \int \log\bigl(f(y|\theta)\bigr) g(z)\,dz \\
    &= -d_{\mathrm{KL}}\bigl(g(z)\|f(z|y,\theta)\bigr) + \log f(y|\theta).
\end{align*}
In other words,
\begin{equation}
  \log f(y|\theta) = \ELBO(g, \theta) + d_{\mathrm{KL}}\bigl(g(z)\|f(z|y,\theta)\bigr).
  \label{eq:10.3}
\end{equation}
The above derivations are direct consequences of Theorem~\ref{thm:10.2} and
Corollary~\ref{cor:10.3} with target $f(z) \equiv f(z|y,\theta)$, unnormalized density
$\tilde{f}(z) \equiv f(y,z|\theta)$, and normalizing constant $c \equiv f(y|\theta)$.

Equation~\eqref{eq:10.3} motivates an iterative approach to maximize the log-likelihood
$\log f(y|\theta)$, which is equivalent to maximizing the likelihood $f(y|\theta)$ since the
logarithm is an increasing function. Given the current iterate $\theta_\ell$, we obtain the next
iterate $\theta_{\ell+1}$ in two steps, maximizing in turn the two components of the ELBO; the
approach is summarized in Figure~\ref{fig:10.2}.

\begin{enumerate}
  \item First, we find $g_\ell(z)$ by maximizing $\ELBO(g, \theta_\ell)$ over p.d.f.\ $g$. From
        \eqref{eq:10.3},
        \[
          \log f(y|\theta_\ell) = \ELBO(g, \theta_\ell) + d_{\mathrm{KL}}\bigl(g\|f(z|y,\theta_\ell)\bigr),
        \]
        and it follows that maximizing $\ELBO(g, \theta_\ell)$ or minimizing
        $d_{\mathrm{KL}}\bigl(g\|f(z|y,\theta_\ell)\bigr)$ over $g$ are equivalent. Since the
        KL-divergence is minimized when both arguments agree, we obtain that
        $g_\ell(z) = f(z|y, \theta_\ell)$.

  \item Second, we find $\theta_{\ell+1}$ by maximizing $\ELBO(g_\ell, \theta)$ over $\theta$. Note
        that
        \begin{equation}
          \ELBO(g_\ell, \theta) = \int \log f(y, z|\theta)\, f(z|y, \theta_\ell)\,dz + \text{const},
          \label{eq:10.4}
        \end{equation}
        where const is independent of $\theta$. Hence, the quantity to maximize is the expected value
        of the joint log-density $\log f(y, z|\theta)$ with respect to $g_\ell(z) = f(z|y, \theta_\ell)$.
\end{enumerate}

Combining these two steps gives the EM algorithm, summarized below.

\begin{figure}[htbp]
  \centering
  \includegraphics[width=0.70\textwidth]{figures/fig_10_2}
  \caption{\textit{One iteration of the EM algorithm.}}
  \label{fig:10.2}
\end{figure}

\begin{algorithm}[H]
  \caption{Expectation Maximization (EM)}
  \label{alg:10.2}
  \begin{algorithmic}[1]
    \Require Initialization $\theta_0$, number of iterations $L$.
    \For{$\ell = 0, 1, \ldots, L-1$}\ do the following expectation and maximization steps:
      \State \textbf{E-Step}: Compute
      \[
        \mathbb{E}_{Z \sim f(z|y,\theta_\ell)}\!\bigl[\log f(y, Z|\theta)\bigr]
        = \int \log f(y, z|\theta)\, f(z|y, \theta_\ell)\,dz.
      \]
      \State \textbf{M-Step}: Compute
      \[
        \theta_{\ell+1} = \argmax_{\theta}\,
        \mathbb{E}_{Z \sim f(z|y,\theta_\ell)}\!\bigl[\log f(y, Z|\theta)\bigr].
      \]
    \EndFor
    \Ensure Parameter $\theta^L$.
  \end{algorithmic}
\end{algorithm}

The following result shows that the EM algorithm has the desirable property that the likelihood
function increases monotonically along iterates $\theta_\ell$.

\begin{theorem}[Monotonic Increase of Likelihood along EM Iterates]
  \label{thm:10.9}
  \textit{Let $\{\theta_\ell\}_{\ell=0}^{L-1}$ be the iterates of Algorithm~\ref{alg:10.2}. Then,
  for $0 \leq \ell \leq L-1$, it holds that}
  \begin{equation}
    \log f(y|\theta_\ell) \leq \log f(y|\theta_{\ell+1}).
    \label{eq:10.5}
  \end{equation}
\end{theorem}

\begin{proof}
Let $g_\ell(z) = f(z|y, \theta_\ell)$. Using the log-likelihood characterization in~\eqref{eq:10.3},
it holds that
\begin{align*}
  \log f(y|\theta_{\ell+1})
    &= \ELBO(g_\ell, \theta_{\ell+1}) + d_{\mathrm{KL}}\bigl(g_\ell\|f(z|y,\theta_{\ell+1})\bigr) \\
    &\geq \ELBO(g_\ell, \theta_{\ell+1}) + d_{\mathrm{KL}}\bigl(g_\ell\|f(z|y,\theta_{\ell+1})\bigr) \\
    &\geq \ELBO(g_\ell, \theta_\ell) + d_{\mathrm{KL}}\bigl(g_\ell\|f(z|y,\theta_{\ell+1})\bigr) \\
    &\geq \ELBO(g_\ell, \theta_\ell) + d_{\mathrm{KL}}\bigl(g_\ell\|f(z|y,\theta_\ell)\bigr) \\
    &= \log f(y|\theta_\ell).
\end{align*}
The first inequality follows because $\theta_{\ell+1} = \argmax_\theta \ELBO(g_\ell, \theta)$; the
second inequality follows because $d_{\mathrm{KL}}\bigl(g_\ell\|f(z|y,\theta_\ell)\bigr) = 0$ since
$g_\ell = f(z|y,\theta_\ell)$, while $d_{\mathrm{KL}}\bigl(g_\ell\|f(z|y,\theta_{\ell+1})\bigr) \geq 0$.
\end{proof}

\begin{remark}
  \label{rem:10.10}
  \textit{Under mild additional assumptions, it follows from Theorem~\ref{thm:10.9} that the
  iterates $\theta_\ell$ of the EM algorithm converge, as $\ell \to \infty$, to a local maximizer
  of the likelihood function. It is important to note, however, that the expectation in the E-step
  and the optimization in the M-step are often intractable. Monte Carlo algorithms may be employed
  to approximate the E-step, and optimization algorithms to approximate the M-step. Such
  approximations can cause loss of monotonicity and convergence guarantees.}
\end{remark}

\medskip
\noindent\fbox{\parbox{\dimexpr\linewidth-2\fboxsep-2\fboxrule\relax}{%
  \textbf{Pros and Cons}: The EM algorithm is a workhorse in computational mathematics and
  statistics. It is widely used to compute maximum likelihood estimators in mixture models,
  hierarchical models, hidden Markov models, etc. For some problems, the M step can be solved in
  closed form, in which case the EM algorithm can be an efficient derivative-free optimization
  method. A caveat of the EM algorithm is that the choice of initialization is important, as the
  algorithm converges to local maxima. In addition, for some problems computing the E and M steps
  can be computationally expensive, and the convergence can be slow.
}}
\medskip

\begin{example}[Dempster et al.\ 1977]
  \label{ex:10.11}
  Observations $y = (y_1, y_2, y_3, y_4)$ are gathered from the multinomial distribution
  \[
    \operatorname{Multinomial}\!\left(K;\, \frac{1}{2}+\frac{\theta}{4},\, \frac{1}{4}(1-\theta),\,
    \frac{1}{4}(1-\theta),\, \frac{\theta}{4}\right).
  \]
  Estimation is simplified if $y_1$ is split into two cells $z = (z_1, z_2)$ so we create the
  augmented model
  \[
    (z_1, z_2, y_2, y_3, y_4) \sim \operatorname{Multinomial}\!\left(K;\,
    \frac{1}{2},\, \frac{\theta}{4},\, \frac{1}{4}(1-\theta),\, \frac{1}{4}(1-\theta),\,
    \frac{\theta}{4}\right),
  \]
  with $y_1 = z_1 + z_2$. Here
  \begin{align*}
    f(y, z|\theta) &\propto \theta^{z_2 + y_4}(1-\theta)^{y_2 + y_3}, \\
    f(y|\theta)    &\propto (2+\theta)^{y_1}\, \theta^{y_4}(1-\theta)^{y_2+y_3}.
  \end{align*}
  We have (up to constants)
  \begin{align*}
    \mathbb{E}_{Z \sim f(z|y,\theta_\ell)}\!\bigl[\log f(y, Z|\theta)\bigr]
      &= \mathbb{E}\bigl[(Z_2 + y_4)\log\theta + (y_2 + y_3)\log(1-\theta)\bigr] \\
      &= \left(\frac{\theta_\ell}{2+\theta_\ell}\, y_1 + y_4\right) \log\theta
         + (y_2 + y_3)\log(1-\theta),
  \end{align*}
  where we used that
  \[
    Z_2 | y, \theta_\ell \sim \operatorname{Binomial}\!\left(y_1,\,
    \frac{\tfrac{\theta_\ell}{4}}{\tfrac{1}{2} + \tfrac{\theta_\ell}{4}}\right)
    = \operatorname{Binomial}\!\left(y_1,\, \frac{\theta_\ell}{2 + \theta_\ell}\right).
  \]
  The expectation $\mathbb{E}_{Z \sim f(z|y,\theta_\ell)}\bigl[\log f(y, Z|\theta)\bigr]$ can then
  be maximized over $\theta$ to give
  \[
    \theta_{\ell+1} = \frac{\dfrac{\theta_\ell}{2+\theta_\ell}\,y_1 + y_4}{\dfrac{\theta_\ell}{2+\theta_\ell}\,y_1 + y_2 + y_3 + y_4}.\qquad\square
  \]
\end{example}

\subsection*{Monte Carlo EM}

In practice, computing the E-step may be computationally difficult, and Monte Carlo can be used to
approximate this expectation:
\[
  \mathbb{E}_{Z \sim f(z|y,\theta_\ell)}\!\bigl[\log f(y, Z|\theta)\bigr]
  \approx \frac{1}{N} \sum_{n=1}^{N} \log f(y, Z^{(n)}|\theta),
\]
where
\[
  Z^{(1)}, \ldots, Z^{(N)} \overset{\mathrm{i.i.d.}}{\sim} f(z|y, \theta_\ell).
\]

\begin{example}[Monte Carlo EM for Dempster et al.\ 1977]
  \label{ex:10.12}
  In the setting of Example~\ref{ex:10.11} there is no need to use Monte Carlo EM, since the
  expectation can be computed analytically using that
  $Z_2|y,\theta_\ell \sim \operatorname{Binomial}\!\bigl(y_1, \tfrac{\theta_\ell}{2+\theta_\ell}\bigr)$.
  However, merely for illustration of the technique, we note that the Monte Carlo version would
  define
  \[
    \theta_{\ell+1} = \frac{\bar{z}_N + y_4}{\bar{z}_N + y_2 + y_3 + y_4},
  \]
  where
  \[
    \bar{z}_N = \frac{1}{N}\sum_{n=1}^{N} Z^{(n)},
    \qquad
    Z^{(1)}, \ldots, Z^{(N)} \overset{\mathrm{i.i.d.}}{\sim}
    \operatorname{Binomial}\!\left(y_1,\, \frac{\theta_\ell}{2+\theta_\ell}\right). \qquad\square
  \]
\end{example}
